%% toc: P3 39 39.3
\chapter{Finite-Variable Languages and Pebble Games}
\label{sec:p03-w03-p03-056-pebble}

A first-order formula may use the same variable symbol many times, binding
it again once its previous value is no longer needed. Counting distinct
variable symbols, rather than quantifiers, gives a hierarchy of languages
whose expressive power is measured exactly by a comparison game played with
a fixed number of pairs of pebbles. This chapter supplies the grammar of the
\(k\)-variable languages, both first-order and infinitary, the \(k\)-pebble
game with its placement, reuse and removal moves, and the finite-variable
equivalence theorem in a version for games of a fixed finite length and a
version for games that go on forever. Every proof is in the annex.

\emph{Inputs.} A finite vocabulary, two structures for it, a number \(k\geq1\)
of variables (equivalently, of pebble pairs), and optionally a number \(n\)
of rounds and an initial placement of pebbles.
\emph{Conclusions.} Which formulas with at most \(k\) variables, and of
quantifier rank at most \(n\) when \(n\) is given, can tell the two pebbled
structures apart; in particular, inexpressibility results for whole languages
obtained by exhibiting winning strategies.

The definitions of the languages, of the game, and of \(k\)-variable
equivalence follow Kolaitis's 2003 manuscript of his survey on the
expressive power of logics on finite models
\cite[pp.~3--4, 52--59]{Kolaitis2003ExpressivePowerNotes}; that
manuscript, not the later printed chapter, is the version followed here.
The finite-round form of the game and of the theorem follows the slides of
the Technion course ``Logical Methods in Combinatorics''
\cite[pp.~10--13]{TechnionLMC2009Lecture4}.

\section{The \(k\)-variable languages}
\label{sec:p03-w03-p03-056-languages}

\begin{hbconvention}{def:p03-w03-p03-056-vocabulary}{Vocabularies and partial isomorphisms}
A \emph{vocabulary} is a finite set
\(\sigma=\{R_1,\ldots,R_m,c_1,\ldots,c_s\}\) of relation symbols with
specified positive arities and constant symbols. A \(\sigma\)-structure
\(\mathbf A\) has a nonempty universe \(A\), a relation
\(R_i^{\mathbf A}\subseteq A^{\operatorname{arity}(R_i)}\) for each relation
symbol and an element \(c_j^{\mathbf A}\in A\) for each constant symbol
\cite[p.~3]{Kolaitis2003ExpressivePowerNotes}. Structures may be
infinite unless a statement says otherwise.

A \emph{partial isomorphism} from \(\mathbf A\) to \(\mathbf B\) is an
isomorphism from a substructure of \(\mathbf A\) to a substructure of
\(\mathbf B\); it therefore sends each \(c_j^{\mathbf A}\) to \(c_j^{\mathbf B}\)
\cite[pp.~3--4]{Kolaitis2003ExpressivePowerNotes}. Equivalently, a
partial isomorphism is an injective map \(h\) with
\(\{c_1^{\mathbf A},\ldots,c_s^{\mathbf A}\}\subseteq\operatorname{dom}(h)
\subseteq A\), \(h(c_j^{\mathbf A})=c_j^{\mathbf B}\), and
\(R^{\mathbf A}(a_1,\ldots,a_r)\iff R^{\mathbf B}(h(a_1),\ldots,h(a_r))\)
for every relation symbol \(R\) of arity \(r\) and all
\(a_1,\ldots,a_r\in\operatorname{dom}(h)\). For a constant-free vocabulary
these are exactly the partial \(\tau\)-embeddings of
\cref{def:p03-w03-p03-039-partial-iso}, with \(\tau=\sigma\).
\end{hbconvention}

\begin{hbdefinition}{def:p03-w03-p03-056-grammar}{The languages \(\mathrm{L}^k_{\infty\omega}\), \(\mathrm{FO}^k_{\omega\omega}\) and \(\mathrm{L}^\omega_{\infty\omega}\)}
Fix a vocabulary \(\sigma\) and an integer \(k\geq1\), and put
\([k]=\{1,\ldots,k\}\). The only variables are \(v_1,\ldots,v_k\). The
\emph{terms} are these variables and the constant symbols of \(\sigma\).
The formulas of \(\mathrm{L}^k_{\infty\omega}(\sigma)\) are generated by
the following clauses and no others.
\begin{enumerate}
\item[(G1)] If \(t,t'\) are terms, \(t=t'\) is an atomic formula; if
      \(R\in\sigma\) has arity \(r\) and \(t_1,\ldots,t_r\) are terms,
      \(R(t_1,\ldots,t_r)\) is an atomic formula.
\item[(G2)] If \(\varphi\) is a formula, so is \(\neg\varphi\).
\item[(G3)] If \(\Phi\) is a \emph{set} of formulas, of any cardinality, then
      \(\bigwedge\Phi\) and \(\bigvee\Phi\) are formulas.
\item[(G4)] If \(\varphi\) is a formula and \(i\in[k]\), then
      \(\exists v_i\,\varphi\) and \(\forall v_i\,\varphi\) are formulas.
\end{enumerate}
Clause (G4) may be applied to \(v_i\) inside the scope of an earlier
quantifier on the same \(v_i\); this is \emph{variable reuse}, and an
occurrence of \(v_i\) is bound by the nearest quantifier on \(v_i\) above
it. The binary connectives are abbreviations: \(\varphi\land\psi\) is
\(\bigwedge\{\varphi,\psi\}\), \(\varphi\lor\psi\) is
\(\bigvee\{\varphi,\psi\}\), \(\varphi\to\psi\) is
\(\bigvee\{\neg\varphi,\psi\}\), and \(t_1\neq t_2\) is \(\neg\,t_1=t_2\).

\(\mathrm{FO}^k_{\omega\omega}(\sigma)\) is the subclass of formulas in whose
formation every application of (G3) has a finite set \(\Phi\); these are the
first-order formulas whose variables are among \(v_1,\ldots,v_k\). Finally
\(\mathrm{L}^\omega_{\infty\omega}=\bigcup_{k\geq1}\mathrm{L}^k_{\infty\omega}\)
\cite[Def.~7.1, pp.~52--53]{Kolaitis2003ExpressivePowerNotes}. Since
\(v_1,\ldots,v_k\) are among \(v_1,\ldots,v_{k'}\) for \(k\leq k'\), every
\(\mathrm{L}^k_{\infty\omega}\)-formula is an
\(\mathrm{L}^{k'}_{\infty\omega}\)-formula.
\index{finite-variable logic}\index{L k infinity omega@\(\mathrm{L}^k_{\infty\omega}\)}
\end{hbdefinition}

Kolaitis defines \(\mathrm{L}^k_{\infty\omega}\) as the
\(\mathrm{L}_{\infty\omega}\)-formulas in which at most \(k\) distinct
variables occur, drawn from an unlimited supply; the grammar above fixes the
\(k\) symbols in advance. \Cref{thm:p03-w03-p03-056-renaming} shows that the
two readings define the same expressive power, so nothing is lost by the
fixed supply. The infinitary language of all set-size conjunctions with
unboundedly many variables is \cref{def:p03-w03-p03-039-language}.

\begin{hbdefinition}{def:p03-w03-p03-056-semantics}{Free labels, assignments, satisfaction and quantifier rank}
The set \(\operatorname{fr}(\varphi)\subseteq[k]\) of \emph{free labels} is
defined by recursion: for an atomic formula it is the set of \(i\) such that
\(v_i\) occurs in it; \(\operatorname{fr}(\neg\varphi)=\operatorname{fr}(\varphi)\);
\(\operatorname{fr}(\bigwedge\Phi)=\operatorname{fr}(\bigvee\Phi)=
\bigcup_{\varphi\in\Phi}\operatorname{fr}(\varphi)\); and
\(\operatorname{fr}(\exists v_i\varphi)=\operatorname{fr}(\forall v_i\varphi)=
\operatorname{fr}(\varphi)\setminus\{i\}\). A \emph{sentence} has no free label.

An \emph{assignment} in \(\mathbf A\) is a map \(\alpha:D\to A\) with
\(D\subseteq[k]\); it is \emph{admissible} for \(\varphi\) when
\(\operatorname{fr}(\varphi)\subseteq D\). For \(i\in[k]\) and \(a\in A\),
\(\alpha[i\mapsto a]\) is the assignment on \(D\cup\{i\}\) that sends \(i\)
to \(a\) and agrees with \(\alpha\) elsewhere. A term has value
\(v_i^{\mathbf A}[\alpha]=\alpha(i)\) and \(c^{\mathbf A}[\alpha]=c^{\mathbf A}\).
For admissible \(\alpha\), \(\mathbf A\models\varphi[\alpha]\) is defined by
recursion on formation: atomic formulas are evaluated at the term values;
\(\neg\) is classical; \(\bigwedge\Phi\) holds when every member holds and
\(\bigvee\Phi\) when some member holds; \(\exists v_i\varphi\) holds when
\(\mathbf A\models\varphi[\alpha[i\mapsto a]]\) for some \(a\in A\), and
\(\forall v_i\varphi\) when this holds for every \(a\in A\). Each recursive
call is again admissible, because subformulas of \(\bigwedge\Phi\) and of
\(\neg\varphi\) have no more free labels and
\(\operatorname{fr}(\varphi)\subseteq\operatorname{fr}(\exists v_i\varphi)\cup\{i\}\).
Kolaitis writes \(\mathbf A\models\varphi(v_1/a_1,\ldots,v_r/a_r)\) for the
case \(D=\{1,\ldots,r\}\), \(\alpha(i)=a_i\)
\cite[Def.~7.8, p.~57]{Kolaitis2003ExpressivePowerNotes}.

The \emph{quantifier rank} \(\operatorname{qr}(\varphi)\) is the ordinal
defined by: \(0\) for atomic formulas;
\(\operatorname{qr}(\neg\varphi)=\operatorname{qr}(\varphi)\);
\(\operatorname{qr}(\bigwedge\Phi)=\operatorname{qr}(\bigvee\Phi)=
\sup_{\varphi\in\Phi}\operatorname{qr}(\varphi)\) (and \(0\) for
\(\Phi=\emptyset\)); and
\(\operatorname{qr}(\exists v_i\varphi)=\operatorname{qr}(\forall v_i\varphi)=
\operatorname{qr}(\varphi)+1\). On \(\mathrm{FO}^k_{\omega\omega}\) this is
the usual finite rank of \cref{def:p03-quantifier-rank}; the extension by
suprema to infinite conjunctions is the reading adopted here for ``quantifier
rank \(n\)'' in the infinitary language. Two formulas with free labels in
\(D\) are \emph{equivalent over \(D\)} when they have the same truth value
in every \(\sigma\)-structure under every assignment with domain \(D\).
\index{quantifier rank!infinitary}
\end{hbdefinition}

\begin{hbdefinition}{def:p03-w03-p03-056-equivalence}{Finite-variable equivalence}
Let \(\alpha:D\to A\) and \(\beta:D\to B\) be assignments with the same
domain \(D\subseteq[k]\). Write
\((\mathbf A,\alpha)\equiv^k_{\infty\omega}(\mathbf B,\beta)\) when
\(\mathbf A\models\varphi[\alpha]\iff\mathbf B\models\varphi[\beta]\) for
every \(\varphi\in\mathrm{L}^k_{\infty\omega}\) with
\(\operatorname{fr}(\varphi)\subseteq D\); define
\((\mathbf A,\alpha)\equiv^k_{\omega\omega}(\mathbf B,\beta)\) in the same
way with \(\mathrm{FO}^k_{\omega\omega}\)-formulas; and, for \(n\in\N\),
define \(\equiv^{k,n}_{\infty\omega}\) and \(\equiv^{k,n}_{\omega\omega}\) by
restricting to formulas of quantifier rank at most \(n\). For \(D=\emptyset\)
we write \(\mathbf A\equiv^k_{\infty\omega}\mathbf B\), and so on: the
structures satisfy the same sentences of the language in question. These are
Kolaitis's \(\mathrm{L}^k_{\infty\omega}\)- and
\(\mathrm{FO}^k_{\omega\omega}\)-equivalence
\cite[Def.~7.8, p.~57]{Kolaitis2003ExpressivePowerNotes}; the slides
write \(\sim^{\infty,\omega}_{k}\) and \(\sim^{\infty,\omega}_{k,n}\) for the
sentence versions without and with the rank bound
\cite[p.~13]{TechnionLMC2009Lecture4}.
\end{hbdefinition}

\begin{hblemma}{thm:p03-w03-p03-056-coincidence}{proved=annex:p03-w03-p03-056-coincidence}{Only free labels matter}
If \(\alpha\) and \(\alpha'\) are assignments in \(\mathbf A\), both
admissible for \(\varphi\in\mathrm{L}^k_{\infty\omega}\), and
\(\alpha(i)=\alpha'(i)\) for every \(i\in\operatorname{fr}(\varphi)\), then
\(\mathbf A\models\varphi[\alpha]\iff\mathbf A\models\varphi[\alpha']\).
\end{hblemma}

\begin{hblemma}{thm:p03-w03-p03-056-renaming}{proved=annex:p03-w03-p03-056-renaming}{Counting variables versus fixing them}
Let \(\varphi\) be a formula of \(\mathrm{L}_{\infty\omega}(\sigma)\), in the
sense of \cref{def:p03-w03-p03-039-language}, in which the set \(W\) of
variables occurring (free or bound) has at most \(k\) members, and let
\(\pi:W\to\{v_1,\ldots,v_k\}\) be injective. Let \(\varphi^\pi\) be the
result of replacing every occurrence of each \(w\in W\) by \(\pi(w)\). Then
\(\varphi^\pi\in\mathrm{L}^k_{\infty\omega}(\sigma)\), and for every
assignment \(s\) of elements of \(A\) to a set \(S\subseteq W\) containing
the free variables of \(\varphi\),
\(\mathbf A\models\varphi[s]\iff\mathbf A\models\varphi^\pi[s\circ\pi^{-1}]\),
where \(s\circ\pi^{-1}\) sends \(\pi(w)\) to \(s(w)\) for \(w\in S\).
\end{hblemma}

\begin{hblemma}{thm:p03-w03-p03-056-walks}{proved=annex:p03-w03-p03-056-walks}{Walks of every length with three variables}
Let \(\sigma=\{E\}\) with \(E\) binary. Put
\(\theta_1=E(v_1,v_2)\) and
\[
 \theta_{n+1}=\exists v_3\bigl(E(v_1,v_3)\land
   \exists v_1\,(v_3=v_1\land\theta_n)\bigr)\qquad(n\geq1).
\]
Each \(\theta_n\) is an \(\mathrm{FO}^3_{\omega\omega}\)-formula with
\(\operatorname{fr}(\theta_n)=\{1,2\}\) and
\(\operatorname{qr}(\theta_n)=2(n-1)\), and
\(\mathbf A\models\theta_n[\alpha]\) holds exactly when there are
\(e_0,\ldots,e_n\in A\) with \(e_0=\alpha(1)\), \(e_n=\alpha(2)\) and
\((e_{j-1},e_j)\in E^{\mathbf A}\) for \(1\leq j\leq n\). Consequently the
\(\mathrm{L}^3_{\infty\omega}\)-sentence
\[
 \forall v_1\forall v_2\Bigl(v_1=v_2\lor\bigvee\{\theta_n:n\geq1\}\Bigr)
\]
holds in a directed graph exactly when every vertex reaches every other one
along \(E\)-edges, which for a symmetric \(E\) is connectivity.
\end{hblemma}

\begin{hbexample}{def:p03-w03-p03-056-reuse-example}{Reuse in the walk formulas}
The first-order formula ``there is a walk of length \(n\) from \(x\) to
\(y\)'' is first written with \(n-1\) existentially quantified intermediate
variables. In \(\theta_{n+1}\) of \cref{thm:p03-w03-p03-056-walks} the
variable \(v_3\) names the first intermediate vertex, the inner
\(\exists v_1\) rebinds \(v_1\) to that same vertex, and the old value of
\(v_1\) is no longer needed; \(\theta_n\) then reuses \(v_1\) and \(v_3\)
for the remaining walk. This is Kolaitis's illustration of reuse
\cite[pp.~53--54]{Kolaitis2003ExpressivePowerNotes}; the formula
displayed above differs from his only in naming \(x,y,z\) as
\(v_1,v_2,v_3\). The quantifier rank grows with \(n\) while the number of
variables stays three, so the infinite disjunction is a three-variable
formula of infinite quantifier rank \(\omega\).
\end{hbexample}

\begin{hblemma}{thm:p03-w03-p03-056-finitely-many}{proved=annex:p03-w03-p03-056-finitely-many}{Finitely many first-order types of bounded rank}
For a vocabulary \(\sigma\), \(k\geq1\), \(m\in\N\) and \(D\subseteq[k]\)
there is a finite set \(F_{m,D}\) of
\(\mathrm{FO}^k_{\omega\omega}(\sigma)\)-formulas, each of quantifier rank
at most \(m\) with free labels in \(D\), such that every
\(\mathrm{FO}^k_{\omega\omega}(\sigma)\)-formula of quantifier rank at most
\(m\) with free labels in \(D\) is equivalent over \(D\) to a member of
\(F_{m,D}\).
\end{hblemma}

The finiteness of \(\sigma\) is used here: the atomic formulas with free
labels in \(D\) form a finite set. With infinitely many relation symbols the
lemma fails already for \(m=0\).

\section{The \(k\)-pebble game}
\label{sec:p03-w03-p03-056-game}

\begin{hbdefinition}{def:p03-w03-p03-056-pebble-game}{The \(k\)-pebble games \(\mathrm{PG}^n_k\) and \(\mathrm{PG}^\infty_k\)}
Let \(\mathbf A,\mathbf B\) be \(\sigma\)-structures and \(k\geq1\). There
are two players, the Spoiler and the Duplicator, and \(k\) pairs of pebbles,
the pair with label \(i\in[k]\) consisting of one pebble for \(\mathbf A\)
and one for \(\mathbf B\).

\emph{Positions.} A position is a pair \((\alpha,\beta)\) of maps
\(\alpha:D\to A\), \(\beta:D\to B\) with a common domain \(D\subseteq[k]\):
the labels in \(D\) are on the board, the \(\mathbf A\)-pebble with label
\(i\) lies on \(\alpha(i)\) and the \(\mathbf B\)-pebble on \(\beta(i)\).
Two labels may lie on the same element.

\emph{Moves.} A round consists of a move of the Spoiler and the answer of
the Duplicator.
\begin{enumerate}
\item[(P)] \emph{Placement.} The Spoiler chooses a label \(i\in[k]\), one of
      the two structures and an element of it, say \(a\in A\). If label
      \(i\) is on the board, both of its pebbles are first lifted; the
      \(\mathbf A\)-pebble \(i\) is put on \(a\), and the Duplicator puts
      the \(\mathbf B\)-pebble \(i\) on an element \(b\in B\) of her choice.
      The new position is \((\alpha[i\mapsto a],\beta[i\mapsto b])\).
      Placement in \(\mathbf B\) is symmetric. Moving a pebble that is
      already in use in this way is \emph{reuse}.
\item[(R)] \emph{Removal.} The Spoiler chooses a label \(i\in D\) and removes
      one of its pebbles; the Duplicator removes the other. The new
      position is \((\alpha\restriction(D\setminus\{i\}),
      \beta\restriction(D\setminus\{i\}))\).
\end{enumerate}

\emph{Winning.} With a position \((\alpha,\beta)\) associate
\[
 \rho_{\alpha,\beta}=\{(\alpha(i),\beta(i)):i\in D\}\cup
 \{(c^{\mathbf A},c^{\mathbf B}):c\in\sigma\text{ a constant symbol}\}.
\]
The position is \emph{good} when \(\rho_{\alpha,\beta}\) is (the graph of) a
partial isomorphism from \(\mathbf A\) to \(\mathbf B\). In
\(\mathrm{PG}^n_k(\mathbf A,\mathbf B)\) from a given position, \(n\in\N\)
rounds are played; in \(\mathrm{PG}^\infty_k(\mathbf A,\mathbf B)\) the play
goes on forever. The Duplicator wins a play when every position in it,
including the initial one and the one after every move, is good; otherwise
the Spoiler wins it. Unless another initial position is named, play starts
from the empty position (\(D=\emptyset\)).

\emph{Strategies.} A strategy for the Duplicator assigns her answer to every
finite sequence of Spoiler moves; it is \emph{winning} from a position when
every play from that position in which she follows it is won by her. If
\(\Sigma\) is winning from \(p\) in a game with \(n\geq1\) rounds (or
forever), and the first round leads, according to \(\Sigma\), to the
position \(p'\), then the strategy
\(h\mapsto\Sigma(\mu\,h)\), where \(\mu\) is that first Spoiler move, is
winning from \(p'\) for \(n-1\) rounds (or forever), since its plays are the
tails of plays of \(\Sigma\).

\emph{Variants.} In the slides the Spoiler only places pebbles, any pebble
being movable, and there are no removal moves
\cite[pp.~10--11]{TechnionLMC2009Lecture4}. In Kolaitis's game a placement
move must use a pebble not currently on the board, and removals are
separate moves \cite[Def.~7.3, p.~55]{Kolaitis2003ExpressivePowerNotes}.
\Cref{thm:p03-w03-p03-056-finite-round,thm:p03-w03-p03-056-eternal} show
that these variants have the same winner as the game above.
\index{pebble game}\index{k-pebble game@\(k\)-pebble game}
\end{hbdefinition}

The Duplicator's goal is the one stated by Kolaitis: to keep the pebbled
correspondence a partial isomorphism at every point of the play
\cite[Def.~7.3, p.~55]{Kolaitis2003ExpressivePowerNotes}. He makes
``winning strategy'' precise as a nonempty family of partial isomorphisms
with at most \(k\) pairs outside the constants, closed under subfunctions
and with the forth and back properties up to \(k\)
\cite[Def.~7.4, p.~55]{Kolaitis2003ExpressivePowerNotes}. The set of
positions used in the proof of \cref{thm:p03-w03-p03-056-eternal} is a family
of that kind, presented by labelled positions.

\begin{hblemma}{thm:p03-w03-p03-056-atomic}{proved=annex:p03-w03-p03-056-atomic}{Good positions are atomic agreement}
A position \((\alpha,\beta)\) with domain \(D\) is good if and only if
\(\mathbf A\models\varphi[\alpha]\iff\mathbf B\models\varphi[\beta]\) for
every atomic formula \(\varphi\) of \(\mathrm{L}^k_{\infty\omega}(\sigma)\)
with \(\operatorname{fr}(\varphi)\subseteq D\).
\end{hblemma}

\begin{hbmethod}{def:p03-w03-p03-056-method}{Proving inexpressibility with pebble games}
\emph{Goal.} Show that a class \(Q\) of \(\sigma\)-structures, closed under
isomorphism, is not the class of models of any
\(\mathrm{L}^\omega_{\infty\omega}\)-sentence within a given class
\(\mathcal C\) of \(\sigma\)-structures.
\emph{Procedure.} For each \(k\geq1\) exhibit
\(\mathbf A_k\in\mathcal C\cap Q\) and \(\mathbf B_k\in\mathcal C\setminus Q\)
and a winning strategy for the Duplicator in
\(\mathrm{PG}^\infty_k(\mathbf A_k,\mathbf B_k)\).
\emph{Side conditions.} The strategy must be winning against every Spoiler
play, including reuse and removals, and must be given for every \(k\); a
strategy for the game of fixed length \(n\) yields only
\(\equiv^{k,n}_{\infty\omega}\) and refutes only sentences of rank at most
\(n\). For a single \(k\) the procedure shows that no
\(\mathrm{L}^k_{\infty\omega}\)-sentence defines \(Q\) on \(\mathcal C\).
\emph{Justification.} A defining sentence would lie in some
\(\mathrm{L}^k_{\infty\omega}\), and \cref{thm:p03-w03-p03-056-eternal} makes
\(\mathbf A_k\) and \(\mathbf B_k\) agree on it. This is the soundness half
of Kolaitis's method of \(k\)-pebble games
\cite[Method~7.11, pp.~59--60]{Kolaitis2003ExpressivePowerNotes}; it
needs no finiteness. When \(\mathcal C\) consists of finite structures the
method is also complete: if \(Q\) is not definable on \(\mathcal C\) by any
\(\mathrm{L}^\omega_{\infty\omega}\)-sentence, then pairs
\(\mathbf A_k,\mathbf B_k\) as in the procedure exist for every \(k\)
(\cref{thm:p03-w03-p03-056-definability}).
\end{hbmethod}

\section{The finite-variable equivalence theorem}
\label{sec:p03-w03-p03-056-theorem}

\begin{hbtheorem}{thm:p03-w03-p03-056-finite-round}{proved=annex:p03-w03-p03-056-finite-round}{Finite-round pebble games and bounded-rank \(k\)-variable formulas}
Let \(\sigma\) be a vocabulary, \(\mathbf A\) and \(\mathbf B\)
\(\sigma\)-structures of arbitrary cardinality, \(k\geq1\), \(n\in\N\), and
\((\alpha,\beta)\) a position with domain \(D\subseteq[k]\). The following
are equivalent.
\begin{enumerate}
\item The Duplicator has a winning strategy in
      \(\mathrm{PG}^n_k(\mathbf A,\mathbf B)\) from \((\alpha,\beta)\).
\item The Duplicator has a winning strategy in the \(n\)-round game from
      \((\alpha,\beta)\) in which the Spoiler makes no removal moves.
\item \((\mathbf A,\alpha)\equiv^{k,n}_{\infty\omega}(\mathbf B,\beta)\).
\item \((\mathbf A,\alpha)\equiv^{k,n}_{\omega\omega}(\mathbf B,\beta)\).
\end{enumerate}
\end{hbtheorem}

The sentence case of the equivalence of items~2 and~3 is Theorem~2(i) of the
slides, stated there with the proof deferred
\cite[p.~13]{TechnionLMC2009Lecture4}; the slides' clause speaks of
sentences ``of quantifier rank \(n\)'', which is read here as rank at most
\(n\). Item~4 adds that, for a fixed number of rounds, infinitary
conjunctions add no distinguishing power.

\begin{hbtheorem}{thm:p03-w03-p03-056-eternal}{proved=annex:p03-w03-p03-056-eternal}{Eternal pebble games and \(\mathrm{L}^k_{\infty\omega}\)-equivalence}
Let \(\sigma\), \(\mathbf A\), \(\mathbf B\), \(k\) and the position
\((\alpha,\beta)\) with domain \(D\) be as in
\cref{thm:p03-w03-p03-056-finite-round}. The following are equivalent.
\begin{enumerate}
\item The Duplicator has a winning strategy in
      \(\mathrm{PG}^\infty_k(\mathbf A,\mathbf B)\) from \((\alpha,\beta)\).
\item She has one in the eternal game without removal moves.
\item She has one in the eternal game in which only pebbles not on the
      board may be placed, removals being separate moves.
\item \((\mathbf A,\alpha)\equiv^k_{\infty\omega}(\mathbf B,\beta)\).
\end{enumerate}
These imply \((\mathbf A,\alpha)\equiv^k_{\omega\omega}(\mathbf B,\beta)\).
If \(\mathbf A\) and \(\mathbf B\) are finite, then
\((\mathbf A,\alpha)\equiv^k_{\omega\omega}(\mathbf B,\beta)\) is equivalent
to items~1--4.
\end{hbtheorem}

For \(D=\emptyset\) this is Kolaitis's Theorem~7.9, which he attributes to
Barwise and to Immerman and proves in outline only
\cite[Thm.~7.9, pp.~58--59]{Kolaitis2003ExpressivePowerNotes}; the
sentence case of items~2 and~4 is Theorem~2(ii) of the slides
\cite[p.~13]{TechnionLMC2009Lecture4}. The finiteness hypothesis in the last
clause cannot be dropped (\cref{thm:p03-w03-p03-056-successor}).

\begin{hbcorollary}{thm:p03-w03-p03-056-monotone}{proved=annex:p03-w03-p03-056-monotone}{Pebble numbers and the equivalence relation}
Let \(1\leq k\leq k'\).
\begin{enumerate}
\item If the Duplicator wins \(\mathrm{PG}^\infty_{k'}(\mathbf A,\mathbf B)\)
      from the empty position, she wins
      \(\mathrm{PG}^\infty_k(\mathbf A,\mathbf B)\) from the empty position;
      equivalently, if the Spoiler wins the \(k\)-pebble game, he wins the
      \(k'\)-pebble game. The same holds for
      \(\mathrm{PG}^n_{k'}\) and \(\mathrm{PG}^n_k\) for each \(n\in\N\).
\item The relation ``the Duplicator wins
      \(\mathrm{PG}^\infty_k(\mathbf A,\mathbf B)\)'' is reflexive, symmetric
      and transitive on \(\sigma\)-structures, and so is the relation ``the
      Duplicator wins \(\mathrm{PG}^n_k(\mathbf A,\mathbf B)\)'' for each
      \(n\in\N\).
\end{enumerate}
\end{hbcorollary}

Kolaitis records both facts as following from the definitions
\cite[p.~56]{Kolaitis2003ExpressivePowerNotes}; here they are read off
the two equivalence theorems, where ``winning'' is the existence of a winning
strategy. In the first item ``the Spoiler wins'' means that the Duplicator
has no winning strategy. A distinguishing infinitary formula then gives the
Spoiler a strategy that reaches a bad position after finitely many rounds
in every play, as the annex explains. On infinite structures there need not
be one finite bound on those round counts that works against every reply;
\cref{thm:p03-w03-p03-056-successor} exhibits this distinction.

\begin{hbtheorem}{thm:p03-w03-p03-056-definability}{proved=annex:p03-w03-p03-056-definability}{Finite-variable definability on classes of finite structures}
Let \(\mathcal C\) be a class of finite \(\sigma\)-structures and \(Q\) a
subclass of \(\mathcal C\) closed under isomorphism within \(\mathcal C\).
The following are equivalent.
\begin{enumerate}
\item Some \(\mathrm{L}^\omega_{\infty\omega}\)-sentence \(\theta\) defines
      \(Q\) on \(\mathcal C\): for \(\mathbf A\in\mathcal C\),
      \(\mathbf A\models\theta\) if and only if \(\mathbf A\in Q\).
\item There is \(k\geq1\) such that for all \(\mathbf A,\mathbf B\in\mathcal C\),
      if \(\mathbf A\in Q\) and the Duplicator wins
      \(\mathrm{PG}^\infty_k(\mathbf A,\mathbf B)\), then \(\mathbf B\in Q\).
\end{enumerate}
When item~2 holds for \(k\), the sentence in item~1 can be taken in
\(\mathrm{L}^k_{\infty\omega}\), as a disjunction of conjunctions of
\(\mathrm{FO}^k_{\omega\omega}\)-sentences.
\end{hbtheorem}

Kolaitis states and proves this as Corollary~7.10
\cite[Cor.~7.10, p.~59]{Kolaitis2003ExpressivePowerNotes}. Where he
forms a disjunction over a list of representatives of the isomorphism types
in \(\mathcal C\), the annex forms it over the set of complete
\(\mathrm{FO}^k_{\omega\omega}\)-theories of members of \(Q\), which avoids
choosing representatives. The implication from item~1 to item~2 does not use
finiteness; the converse does, through the last clause of
\cref{thm:p03-w03-p03-056-eternal}.

\section{Worked examples}
\label{sec:p03-w03-p03-056-examples}

For \(m\geq1\), \(\mathbf K_m\) is the complete graph on \(m\) vertices: a
\(\{E\}\)-structure with \(m\) elements in which \(E\) holds exactly between
distinct elements.

\begin{hbtheorem}{thm:p03-w03-p03-056-cliques}{proved=annex:p03-w03-p03-056-cliques}{Cliques and pebble numbers}
For every \(k\geq1\):
\begin{enumerate}
\item the Duplicator wins \(\mathrm{PG}^\infty_k(\mathbf K_k,\mathbf K_{k+1})\),
      so \(\mathbf K_k\equiv^k_{\infty\omega}\mathbf K_{k+1}\);
\item the Duplicator has no winning strategy in
      \(\mathrm{PG}^{k+1}_{k+1}(\mathbf K_k,\mathbf K_{k+1})\), and the
      \(\mathrm{FO}^{k+1}_{\omega\omega}\)-sentence
      \(\exists v_1\cdots\exists v_{k+1}\bigwedge\{v_i\neq v_j:1\leq i<j\leq k+1\}\),
      of quantifier rank \(k+1\), holds in \(\mathbf K_{k+1}\) and fails in
      \(\mathbf K_k\).
\end{enumerate}
\end{hbtheorem}

These are the two assertions of Kolaitis's Example~7.5, which he calls
``quite clear'' \cite[Ex.~7.5, p.~56]{Kolaitis2003ExpressivePowerNotes};
the annex gives the strategies.

\begin{hbcorollary}{thm:p03-w03-p03-056-even}{proved=annex:p03-w03-p03-056-even}{Even size is not finite-variable definable}
There is no \(\mathrm{L}^\omega_{\infty\omega}\)-sentence \(\theta\) over
\(\{E\}\) such that, for every finite graph \(\mathbf G\),
\(\mathbf G\models\theta\) exactly when the universe of \(\mathbf G\) has
even size.
\end{hbcorollary}

This is the first item of Kolaitis's Proposition~7.12, where it is combined
with the inclusion of least and partial fixed-point logic in
\(\mathrm{L}^\omega_{\infty\omega}\) on finite structures to conclude that
neither fixed-point logic expresses even cardinality
\cite[Thm.~7.2, p.~54; Prop.~7.12, p.~60]{Kolaitis2003ExpressivePowerNotes}.
That inclusion is quoted, not proved, here.

\begin{hbcorollary}{thm:p03-w03-p03-056-clique-query}{proved=annex:p03-w03-p03-056-clique-query}{Clique queries need many variables}
For every \(k\geq1\) there is no \(\mathrm{L}^k_{\infty\omega}\)-sentence
\(\theta\) over \(\{E\}\) such that, for every finite graph \(\mathbf G\)
(an \(\{E\}\)-structure with \(E\) symmetric and irreflexive),
\(\mathbf G\models\theta\) exactly when \(\mathbf G\) contains \(k+1\)
pairwise adjacent vertices. On graphs, where \(E\) is symmetric, the query is defined by the
\(\mathrm{FO}^{k+1}_{\omega\omega}\)-sentence
\(\exists v_1\cdots\exists v_{k+1}\bigwedge\{E(v_i,v_j):1\leq i<j\leq k+1\}\).
\end{hbcorollary}

This is the second item of Kolaitis's Proposition~7.12
\cite[Prop.~7.12, p.~60]{Kolaitis2003ExpressivePowerNotes}.

For \(m\geq1\), \(\mathbf L_m\) is the linear order \((\{1,\ldots,m\},<)\).

\begin{hbtheorem}{thm:p03-w03-p03-056-orders}{proved=annex:p03-w03-p03-056-orders}{Two pebbles count the length of a finite linear order}
If \(1\leq m<n\), the Duplicator has no winning strategy in
\(\mathrm{PG}^n_2(\mathbf L_m,\mathbf L_n)\). Hence some
\(\mathrm{FO}^2_{\omega\omega}\)-sentence of quantifier rank at most \(n\)
holds in exactly one of \(\mathbf L_m\) and \(\mathbf L_n\).
\end{hbtheorem}

Kolaitis describes the Spoiler's climbing strategy and contrasts it with the
Ehrenfeucht--Fra\"iss\'e game, in which the Duplicator wins the \(r\)-move
game on sufficiently long finite linear orders
\cite[Ex.~7.6, p.~56]{Kolaitis2003ExpressivePowerNotes}. Two variables
reused \(n\) times thus distinguish what \(n\) fresh variables do not
distinguish when the orders are long. He also notes that on finite linear
orders the sentence ``there are at least \(m\) elements'' is equivalent to
an \(\mathrm{L}^2_{\infty\omega}\)-sentence, so that every set of
cardinalities is \(\mathrm{L}^2_{\infty\omega}\)-definable on finite linear
orders \cite[p.~53]{Kolaitis2003ExpressivePowerNotes}.

\section{Limitations}
\label{sec:p03-w03-p03-056-limits}

\begin{hbtheorem}{thm:p03-w03-p03-056-successor}{proved=annex:p03-w03-p03-056-successor}{First-order two-variable equivalence without infinitary equivalence}
Let \(\sigma=\{E\}\). Let \(\mathbf A\) have universe \(\N\) and
\(E^{\mathbf A}=\{(x,x+1):x\in\N\}\), and let \(\mathbf B\) have universe
\((\{0\}\times\N)\cup(\{1\}\times\mathbb Z)\) and
\(E^{\mathbf B}=\{((e,x),(e,x+1))\}\). Then
\begin{enumerate}
\item for every \(n\in\N\) the Duplicator wins
      \(\mathrm{PG}^n_2(\mathbf A,\mathbf B)\), so
      \(\mathbf A\equiv^2_{\omega\omega}\mathbf B\);
\item with \(\delta_0=(v_1=v_1)\) and
      \(\delta_{j+1}=\exists v_2\bigl(E(v_2,v_1)\land
      \exists v_1(v_1=v_2\land\delta_j)\bigr)\), the
      \(\mathrm{L}^2_{\infty\omega}\)-sentence
      \(\exists v_1\bigwedge\{\delta_j:j\in\N\}\) holds in \(\mathbf B\) and
      fails in \(\mathbf A\); so
      \(\mathbf A\not\equiv^2_{\infty\omega}\mathbf B\) and the Duplicator
      has no winning strategy in \(\mathrm{PG}^\infty_2(\mathbf A,\mathbf B)\).
\end{enumerate}
\end{hbtheorem}

\begin{hbwarning}{def:p03-w03-p03-056-limits}{What the pebble game does not measure}
\begin{enumerate}
\item \emph{Finite games and the eternal game differ on infinite
      structures.} \Cref{thm:p03-w03-p03-056-successor} shows that winning
      every finite-length game, which is
      \(\equiv^k_{\omega\omega}\), does not give a win in the eternal game.
      The last clause of \cref{thm:p03-w03-p03-056-eternal} needs both
      structures finite: its proof forms a conjunction indexed by the
      elements of one structure.
\item \emph{Variables, not quantifier rank.} An
      \(\mathrm{FO}^k_{\omega\omega}\)-formula may have any quantifier
      rank, and an \(\mathrm{L}^k_{\infty\omega}\)-formula infinite rank
      \cite[p.~53]{Kolaitis2003ExpressivePowerNotes}. The eternal game
      controls only the number of variables. On finite structures,
      \(\mathbf L_m\) and \(\mathbf L_n\) for \(m<n\) are separated by two
      variables (\cref{thm:p03-w03-p03-056-orders}).
\item \emph{Counting.} \(\mathbf K_k\) and \(\mathbf K_{k+1}\) are
      \(\mathrm{L}^k_{\infty\omega}\)-equivalent
      (\cref{thm:p03-w03-p03-056-cliques}), yet differ in size. Gr\"adel's
      \(C^k_{\infty\omega}\) adds to \(\mathrm{L}^k_{\infty\omega}\) the
      quantifiers \(\exists^{\geq m}\) (``there exist at least \(m\)'')
      \cite[p.~190]{Gradel2007FMTDescriptiveChapter}; the one-variable
      sentence \(\exists^{\geq k+1}v_1\,(v_1=v_1)\) separates the two
      cliques. The \(k\)-pebble game is therefore not a game for the
      counting logics.
\item \emph{Vocabulary.} The vocabulary is finite and has no function
      symbols. \Cref{thm:p03-w03-p03-056-finitely-many}, and with it the
      equivalence of items~3 and~4 of
      \cref{thm:p03-w03-p03-056-finite-round}, uses that finiteness.
\end{enumerate}
\end{hbwarning}

\section{Relationships}
\label{sec:p03-w03-p03-056-relations}

\begin{hblemma}{thm:p03-w03-p03-056-ef-bridge}{proved=annex:p03-w03-p03-056-ef-bridge}{Pebble and Ehrenfeucht--Fra\"iss\'e games}
For a constant-free finite relational vocabulary, \(\sigma\)-structures
\(\mathbf A,\mathbf B\) and \(n\geq1\), the Duplicator has a winning
strategy in \(\mathrm{PG}^n_n(\mathbf A,\mathbf B)\) if and only if she has
one in the \(n\)-round Ehrenfeucht--Fra\"iss\'e game of
\cref{def:p03-ef-game}.
\end{hblemma}

The slides record the same identification
\cite[p.~11]{TechnionLMC2009Lecture4}. Combined with
\cref{thm:p03-w03-p03-056-finite-round}, agreement through rank \(n\) in the
sense of \cref{def:p03-quantifier-rank} and \cref{thm:p03-ef} becomes
agreement on \(n\)-variable formulas of rank at most \(n\).

\begin{itemize}
\item \emph{Infinitary logic.} \(\mathrm{L}^k_{\infty\omega}\) is a fragment of
      \(\mathrm{L}_{\infty\omega}\) (\cref{def:p03-w03-p03-039-language},
      via \cref{thm:p03-w03-p03-056-renaming}). Hence structures that are
      \(\equiv_{\infty,\omega}\)-equivalent, characterized by partial
      isomorphism systems in Karp's theorem
      (\cref{thm:p03-w03-p03-039-karp}), are
      \(\equiv^k_{\infty\omega}\)-equivalent for every \(k\). The
      \(k\)-pebble game is the restriction of Karp's game in which only
      \(k\) pairs may be held at once.
\item \emph{Two-variable logic.} The sentences of
      \(\mathrm{FO}^2_{\omega\omega}\) over a constant-free relational
      vocabulary are the \(\mathrm{FO}^2\)-sentences of
      \cref{def:p02-w02-p02-035-fo2}; \cref{thm:p03-w03-p03-056-finite-round}
      with \(k=2\) is therefore a game criterion for that logic.
\item \emph{Width.} Gr\"adel measures a first-order formula by its width, the
      largest number of free variables of a subformula, and shows that width
      \(k\) is the same as being equivalent, by renaming bound variables, to
      a formula with at most \(k\) distinct variables
      \cite[Lemma~3.1.4, p.~133]{Gradel2007FMTDescriptiveChapter}.
\item \emph{Zero--one laws and fixed points.} The route to the zero--one law
      for \(\mathrm{L}^\omega_{\infty\omega}\) recorded in
      \cref{sec:p09-presentation-n-c-002} uses Duplicator strategies in the
      \(k\)-pebble game on structures satisfying extension axioms; the
      game theorem it relies on is
      \cref{thm:p03-w03-p03-056-eternal}. Kolaitis states that every stage
      of a first-order operator defined with \(k\) variables is
      \(\mathrm{FO}^k_{\omega\omega}\)-definable and that partial fixed
      points are \(\mathrm{L}^k_{\infty\omega}\)-definable on finite
      structures, referring elsewhere for the proof
      \cite[Thm.~7.2, p.~54]{Kolaitis2003ExpressivePowerNotes}.
\end{itemize}

\paragraph{Sources.}
Languages, game, equivalence and examples:
\cite[pp.~3--4, 52--60]{Kolaitis2003ExpressivePowerNotes}, a manuscript
dated 1 August 2003. Finite-length games and the finite-round statement:
\cite[pp.~10--13]{TechnionLMC2009Lecture4} and
\cite[p.~2]{TechnionLMC2009Lecture5}. Width and counting logics:
\cite[pp.~133, 190]{Gradel2007FMTDescriptiveChapter}. The proofs in the annex
are written out for this handbook; Kolaitis's proof of the eternal theorem is
an outline, and the slides defer theirs.
